\textbf{Phần 1. Trắc nghiệm lựa chọn}
\setcounterpageref{ex}{0}

\begin{ex} %[4H6N32]
Hai đường thẳng vuông góc với nhau tạo thành
\choice
{Góc nhọn}
{Góc tù}
{Góc bẹt}
{\True Góc vuông ($90^\circ$)}
\loigiai{
Tạo thành góc vuông $90^\circ$.\\
\textbf{Đáp án của câu hỏi là D}
}
\end{ex}

\begin{ex} %[4H6N32]
Hai đường thẳng song song là hai đường thẳng
\choice
{Vuông góc nhau}
{\True Không bao giờ cắt nhau}
{Cắt nhau tại 1 điểm}
{Có 4 góc vuông}
\loigiai{
Hai đường thẳng song song không có điểm chung.\\
\textbf{Đáp án của câu hỏi là B}
}
\end{ex}

\begin{ex} %[4H6H32]
Hình nào sau đây có 2 đường chéo vuông góc với nhau?
\choice
{Hình chữ nhật}
{Hình bình hành}
{Hình thang}
{\True Hình thoi (và hình vuông)}
\loigiai{
Hình thoi có hai đường chéo vuông góc với nhau.\\
\textbf{Đáp án của câu hỏi là D}
}
\end{ex}

\begin{ex} %[4H6H32]
Một hình bình hành có độ dài đáy $20\text{ cm}$, chiều cao $12\text{ cm}$. Diện tích của nó là
\choice
{\True $240\text{ cm}^2$}
{$120\text{ cm}^2$}
{$64\text{ cm}^2$}
{$320\text{ cm}^2$}
\loigiai{
$S = 20 \times 12 = 240\text{ cm}^2$.\\
\textbf{Đáp án của câu hỏi là A}
}
\end{ex}

\begin{ex} %[4H6H32]
Chu vi hình thoi có cạnh dài $16\text{ dm}$ là
\choice
{$32\text{ dm}$}
{\True $64\text{ dm}$}
{$48\text{ dm}$}
{$60\text{ dm}$}
\loigiai{
$P = 16 \times 4 = 64\text{ dm}$.\\
\textbf{Đáp án của câu hỏi là B}
}
\end{ex}

\begin{ex} %[4H6H32]
Trong hình bình hành $ABCD$, cặp cạnh nào sau đây song song với nhau?
\choice
{\True $AB$ và $CD$}
{$AB$ và $BC$}
{$BC$ và $CD$}
{$AD$ và $AB$}
\loigiai{
$AB // CD$.\\
\textbf{Đáp án của câu hỏi là A}
}
\end{ex}

\begin{ex} %[4H6V32]
Một mảnh đất hình bình hành có đáy $35\text{ m}$, chiều cao $20\text{ m}$. Diện tích mảnh đất đó là
\choice
{$350\text{ m}^2$}
{$1400\text{ m}^2$}
{\True $700\text{ m}^2$}
{$110\text{ m}^2$}
\loigiai{
$S = 35 \times 20 = 700\text{ m}^2$.\\
\textbf{Đáp án của câu hỏi là C}
}
\end{ex}

\begin{ex} %[4H6V32]
Một hình thoi có chu vi bằng chu vi hình vuông cạnh $15\text{ cm}$. Độ dài cạnh của hình thoi là
\choice
{\True $15\text{ cm}$}
{$10\text{ cm}$}
{$20\text{ cm}$}
{$30\text{ cm}$}
\loigiai{
Cả 2 hình đều có $P = 4a$ nên cạnh hình thoi cũng bằng $15\text{ cm}$.\\
\textbf{Đáp án của câu hỏi là A}
}
\end{ex}

\begin{ex} %[4H6V32]
Hai đường thẳng $a$ và $b$ cùng vuông góc với đường thẳng $c$. Khi đó $a$ và $b$
\choice
{Vuông góc với nhau}
{Cắt nhau}
{\True Song song với nhau}
{Trùng nhau}
\loigiai{
$a // b$.\\
\textbf{Đáp án của câu hỏi là C}
}
\end{ex}

\begin{ex} %[4H6V32]
Một hình tam giác vuông có hai cạnh góc vuông dài $6\text{ cm}$ và $8\text{ cm}$. Diện tích của hình tam giác vuông đó là
\choice
{$48\text{ cm}^2$}
{$14\text{ cm}^2$}
{$28\text{ cm}^2$}
{\True $24\text{ cm}^2$}
\loigiai{
$S = 6 \times 8 : 2 = 24\text{ cm}^2$.\\
\textbf{Đáp án của câu hỏi là D}
}
\end{ex}

\begin{ex} %[4H6V32]
Số góc vuông trong một hình chữ nhật là
\choice
{$2$}
{\True $4$}
{$3$}
{$1$}
\loigiai{
Có 4 góc vuông.\\
\textbf{Đáp án của câu hỏi là B}
}
\end{ex}

\begin{ex} %[4H6C32]
Một mảnh vườn hình chữ nhật có chu vi $120\text{ m}$, chiều dài hơn chiều rộng $20\text{ m}$. Người ta chia mảnh vườn thành hai phần: một phần hình vuông có cạnh bằng chiều rộng mảnh vườn để trồng rau, phần còn lại trồng hoa. Diện tích phần đất trồng hoa là
\choice
{$300\text{ m}^2$}
{$500\text{ m}^2$}
{\True $400\text{ m}^2$}
{$600\text{ m}^2$}
\loigiai{
Nửa chu vi: $60\text{ m}$. Dài: $40\text{ m}$, Rộng: $20\text{ m}$.\\
Diện tích trồng rau (hình vuông $20 \times 20$): $400\text{ m}^2$.\\
Diện tích cả vườn: $40 \times 20 = 800\text{ m}^2$.\\
Diện tích trồng hoa: $800 - 400 = 400\text{ m}^2$.\\
\textbf{Đáp án của câu hỏi là C}
}
\end{ex}

\textbf{Phần 2. Trắc nghiệm đúng sai}
\setcounterpageref{ex}{0}

\begin{ex} %[4H6H32]
Xét tính đúng sai của các đặc điểm hình học:
\choiceTF
{\True Hình vuông là hình thoi đặc biệt có 4 góc vuông.}
{\True Hình vuông là hình chữ nhật đặc biệt có 4 cạnh bằng nhau.}
{Hình bình hành luôn có 4 cạnh bằng nhau.}
{\True Hình thoi có 2 đường chéo vuông góc với nhau.}
\loigiai{
\begin{itemize}
    \item Hình vuông là hình thoi đặc biệt: Đúng.
    \item Hình vuông là hình chữ nhật đặc biệt: Đúng.
    \item Hình bình hành chỉ có các cạnh đối bằng nhau (không nhất thiết 4 cạnh bằng nhau). Sai.
    \item Hai đường chéo hình thoi vuông góc: Đúng.
\end{itemize}
\begin{itemize}
    \item a) \textbf{Đúng}.
    \item b) \textbf{Đúng}.
    \item c) \textbf{Sai}.
    \item d) \textbf{Đúng}.
\end{itemize}
}
\end{ex}

\begin{ex} %[4H6V32]
Một hình bình hành $ABCD$ có đáy $CD = 16\text{ cm}$, chiều cao $AH = 9\text{ cm}$, cạnh bên $BC = 12\text{ cm}$. Xét tính đúng sai:
\choiceTF
{\True Chu vi hình bình hành là $56\text{ cm}$.}
{\True Diện tích hình bình hành là $144\text{ cm}^2$.}
{Đoạn thẳng $AH$ song song với cạnh đáy $CD$.}
{\True Đoạn thẳng $AH$ vuông góc với cạnh đáy $CD$.}
\loigiai{
\begin{itemize}
    \item Chu vi: $(16 + 12) \times 2 = 56\text{ cm}$. Đúng.
    \item Diện tích: $16 \times 9 = 144\text{ cm}^2$. Đúng.
    \item Đường cao $AH$ vuông góc với đáy $CD$ (không song song). Ý 3 Sai, Ý 4 Đúng.
\end{itemize}
\begin{itemize}
    \item a) \textbf{Đúng}.
    \item b) \textbf{Đúng}.
    \item c) \textbf{Sai}.
    \item d) \textbf{Đúng}.
\end{itemize}
}
\end{ex}

\textbf{Phần 3. Trả lời ngắn}
\setcounterpageref{ex}{0}

\begin{ex} %[4H6H32]
Một hình bình hành có đáy $25\text{ cm}$, chiều cao $14\text{ cm}$. Tính diện tích theo xăng-ti-mét vuông.
\shortans{350}
\loigiai{
\textbf{Đáp số:} $350$\\
$25 \times 14 = 350\text{ cm}^2$.
}
\end{ex}

\begin{ex} %[4H6H32]
Một hình thoi có cạnh $18\text{ cm}$. Tính chu vi hình thoi theo xăng-ti-mét.
\shortans{72}
\loigiai{
\textbf{Đáp số:} $72$\\
$18 \times 4 = 72\text{ cm}$.
}
\end{ex}

\begin{ex} %[4H6V32]
Một mảnh đất hình bình hành có diện tích $450\text{ m}^2$, chiều cao $15\text{ m}$. Tính độ dài đáy theo mét.
\shortans{30}
\loigiai{
\textbf{Đáp số:} $30$\\
$450 : 15 = 30\text{ m}$.
}
\end{ex}

\begin{ex} %[4H6V32]
Một hình thoi có chu vi $100\text{ cm}$. Tính độ dài cạnh hình thoi theo đề-xi-mét (dùng dấu phẩy nếu có).
\shortans{2,5}
\loigiai{
\textbf{Đáp số:} $2{,}5$\\
Cạnh: $100 : 4 = 25\text{ cm} = 2{,}5\text{ dm}$.
}
\end{ex}

\begin{ex} %[4H6V32]
Hình chữ nhật có bao nhiêu cặp cạnh song song với nhau?
\shortans{2}
\loigiai{
\textbf{Đáp số:} $2$\\
2 cặp.
}
\end{ex}

\begin{ex} %[4H6C32]
Một hình vuông có diện tích $64\text{ cm}^2$. Người ta ghép 3 hình vuông như thế thành một hình chữ nhật. Tính chu vi hình chữ nhật đó theo xăng-ti-mét.
\shortans{64}
\loigiai{
\textbf{Đáp số:} $64$\\
Cạnh hình vuông: $8\text{ cm}$.\\
Hình chữ nhật có kích thước $24\text{ cm}$ và $8\text{ cm}$.\\
Chu vi: $(24 + 8) \times 2 = 64\text{ cm}$.
}
\end{ex}
